\begin{table}[H]
\centering
\caption{Diferença entre a proporção de domicílios com internet e com esgotamento adequado, Sergipe, 2022}
\label{tab:item2-7-internet-esgoto}
\small
\begin{tabular}{lrrr}
\toprule
Município & Internet (\%) & Esgotamento adequado (\%) & Diferença (p.p.) \\
\midrule
\multicolumn{4}{l}{\textbf{Maiores diferenças positivas}} \\
Santana do São Francisco & 81,72 & 3,53 & 78,19 \\
Areia Branca & 78,88 & 1,81 & 77,07 \\
Muribeca & 75,72 & 6,13 & 69,59 \\
\midrule
\multicolumn{4}{l}{\textbf{Menores diferenças positivas}} \\
São Miguel do Aleixo & 64,45 & 61,80 & 2,65 \\
Rosário do Catete & 83,64 & 83,03 & 0,61 \\
Carmópolis & 84,52 & 84,30 & 0,22 \\
\midrule
Sergipe & 81,34 & 67,04 & 14,30 \\
\bottomrule
\end{tabular}
\fonte{Elaboração própria com base nos microdados do Censo Demográfico 2022.}
\nota{Esgotamento considerado adequado para valores da variável D0250 até 3, seguindo o critério do exercício, incluindo fossa séptica ou fossa-filtro não ligada à rede (valor 3); diferença em pontos percentuais.}
\end{table}
